\section{Classical boundaries and scope}\label{sec:comparison}
\subsection{Transformation groups}
Gorodski--Lytchak \cite[Section 1 and Proposition 1.1]{GLorbit} distinguish metric quotient-equivalence from ambient orbit-equivalence and study reductions, boundary, cohomogeneity and copolarity. Their Proposition 1.1 imposes a no-boundary hypothesis when recovering the ambient dimension. The half-line example in Proposition~\ref{prop:quotients} has boundary and makes no contrary assertion. Their principal-isotropy reductions and polar examples identify quotient metric spaces; they do not provide command intertwiners preserving specified numerical words. The full alphabet is additional dynamic data.

The sphere-orbifold classification in \cite[Theorem 1.1]{GLorbifold} concerns the geometry of the quotient, not the wordwise positive width of a chosen generating set. Gorodski--Thorbergsson \cite[Theorems 1.1 and 1.2]{GTorbit} classify variationally complete and irreducible taut representations through orbit equivalence and osculating-space structure. Those results can identify useful orbit types in their classes; neither tautness nor polarity supplies a gap for a given alphabet. Our uniform orbital small-ball estimate is an inherited compact differential-chart argument. It is not presented as a replacement for the slice theorem or a new classification of compact representations. The algebraic evaluation of selected adjoint and exterior-power minima in the supporting manuscript remains intact; no general highest-weight classification is asserted here.

The support-net compiler uses the extrinsic convex hull of an orbit and quadratic loss at a support maximum. These are standard convex and Riemannian geometric facts \cite{Barvinok,Orbitopes,Iacobelli}. Its role is to make one legal set of transition rows serve every layer. The nongapped hierarchy is more sharply alphabet-dependent: even preserving the full orbit foliation in the same ambient plane does not preserve its width exponent.

\subsection{Spectral design, circle walks and branching programs}
Boyd--Diaconis--Xiao \cite[Sections 2--4]{BDX} formulate finite reversible Markov-chain design as convex spectral-norm minimization with semidefinite duality. Proposition~\ref{prop:gap-dual} uses that classical convex structure, but its admissible operators act on an infinite spherical harmonic sum. Its truncated programs have the opposite certification direction from an upper tail estimate. We do not infer a positive infinite gap from a finite SDP.

The algebraic expansion theorem of Benoist--de Saxc\'e \cite[Theorem 1.2 and Definition 2.1]{BdS} is an external qualitative input. In Theorem~\ref{thm:binary-example} it is used only for the symmetric average of the explicitly verified dense generators $b,c$. Our four-word variance estimate then supplies the stated two-step norm loss. This separation avoids treating the norm of a nonnormal one-step average as its spectral radius.

Berkes--Borda \cite{BB} analyze circle random walks using Fourier multipliers and Diophantine data. The simultaneous pigeonhole argument and the algebraic-norm separation in Section~\ref{sec:toral} are elementary arithmetic tools. The additional assertion here concerns loss forced by compression into an arbitrary hidden register and a matching exact positive realization. Its test packets are independent as complete words, not independent-letter samples inside a packet. The $r=1$ mechanism and the general conditional-centroid packet inequality are inherited from the earlier rotation development; the multivariate separation, every-scale denominator argument and their matched $r/(2r+1)$ law are the extension proved here.

The clocked model is a nonuniform probabilistic ordered read-once layered program with real terminal columns. Regular and permutation branching-program Fourier bounds, such as \cite{RSV,SVW,LPV}, concern their stated transition assumptions and pseudorandomness or Fourier growth. This paper does not claim those objectives or count total program size. A Bernoulli row converts one bounded column into randomized Boolean acceptance, but preserving a cutpoint sign is weaker than preserving its numerical value. The constant-width attenuated sign simulator, the version-pinned quantum-automata comparison and the projective examples remain in the unaltered supporting manuscript; they are not new claims of this revision.

\subsection{Resource ledger in the theorem}
For a common-row upper realization with $K$ labels there are $|\mathcal A|K$ command rows, in addition to the seed and decoder rows. A bound of $d+1$ successors per row does not bound the complexity of constructing their weights or storing their real values. For a general clocked machine there can instead be $|\mathcal A|\sum_{t<N}K_t$ distinct command rows. None of these descriptions is charged as persistent label width.

The largest toral widths in Theorem~\ref{thm:main} correspond to
$\frac{r}{2r+1}\log_2N+O(1)$ label bits, not polynomially many bits. Exact irrational stochastic rows are not finite random-bit algorithms. Output error, calibration, alphabet size and representation dimension are fixed before $N$ grows. A simultaneous channel of all coordinate answers, repeated queries, an anytime decoder and adaptive experimental design are different tasks.

\subsection{Conclusions}
The alphabet can be expanding only after a finite word block, or nongapped at every finite block while still forcing a sharp growing width. Optimizing the proof law makes this distinction canonical. The explicit nongapped hierarchy shows why a representation exponent cannot classify all alphabets, and supplies matching laws in a substantial excluded family rather than merely restating the exclusion.

The numerical calibration is now a finite algebraic problem with certified convex upper bounds and quantitative frame stability. This does not make spectral-gap positivity decidable. The general dominant-nongapped problem beyond the dual badly approximable planar families, sharp finite integer widths, uniform varying-signal laws and optimal table or fair-bit costs remain distinct questions. The optional arithmetic LPS application in the preserved manuscript retains its original explicitly stated external dependency and its original-page audit limitation. Neither these results nor their reproducible builds assert an independent priority certification or an editorial outcome.
